\begin{abstract}
Existing geometry benchmarks measure whether language models can \emph{interpret} a diagram and answer a question about it; none measures whether a model can \emph{produce} a correct diagram from a description --- the deployed task when LLMs author worksheets, slides, and technical illustrations.
We introduce \textbf{\GeoGen}, a methodology for building automatically-verifiable benchmarks of generative geometric reasoning, and \textbf{\GeoGenBench}, its first instantiation.
\GeoGen pairs \emph{procedural co-generation} of prompts and rubrics from shared source code (eliminating drift between question and grading criteria) with a \emph{render-then-symbolic-verify} pipeline that compiles the model's TikZ output to symbolic objects and decides declarative geometric predicates in milliseconds --- no human grading, no vision-language judge.
\GeoGenBench instantiates the methodology for plane Euclidean geometry: \textbf{801} natural-language prompts paired with machine-checkable properties drawn from a 17-predicate language.
Across six frontier models from three vendors and two generation strategies, \emph{the cost--quality Pareto frontier rewards mid-tier scaffolding over both cheap-model and frontier-model configurations}: under our IR strategy, three frontier models converge within $1$\pp of each other --- \gpt ($94.1$\%, \dollars{0.075}/scenario), \sonnet ($93.4$\%, \dollars{0.088}), and \opus ($93.3$\%, \dollars{0.532}) --- so the cost--quality Pareto frontier collapses to a single corner with \gpt and \sonnet essentially tied at $\sim$$\mathbf{7{\times}}$ lower cost than \opus, and on raw TikZ \gpt at \dollars{0.058}/scenario beats \opus at \dollars{0.262} by $20$\pp ($91.9$\% vs $71.6$\%).
A generic \emph{intermediate-representation} (IR) strategy --- in which the model emits a typed JSON description of the construction (points, lines, circles, intersections) that we compile to TikZ and verify before rendering --- adds $2$--$25$\pp on top of raw TikZ on strict pass rate, with the largest gains on the models that struggle most at \rawcode; one vendor (Gemini Flash) rejects the IR as exceeding its structured-output complexity cap, a deployment-portability gap rather than a capability gap.
We audit the verifier via adversarial probing (identifying a class of \emph{self-referential} predicates that pass tautologically) and a pre-registered $N{=}200$ human-correlation study with two domain-expert raters: the auto-verdict moderately agrees with expert consensus ($\kappa{=}0.455$ on the binary collapse, $95\%$ CI $[0.30, 0.59]$) and is one-directionally lenient ($\sim 18\%$ of \texttt{pass} records overturn to expert \texttt{fail}, zero records the reverse), with the disagreement concentrated on auxiliary-construction templates whose property lists do not enforce prompt-specified numeric parameters.
We release the dataset, the procedural template engine that generated it (allowing reviewers to roll new contamination-free splits), the verification toolkit, and a public leaderboard.
\end{abstract}
